Una partícula $\alpha$ ($^{4}_{2}$He) es dirigeix directament cap al nucli d'un àtom d'urani ($^{238}_{92}$U). El radi del nucli d'urani és, aproximadament, de 0,008~pm (picòmetres).

\figura[0.8]{fig1.png}

\begin{apartat}{1}
Compareu quantitativament els valors del mòdul de la intensitat del camp elèctric degut al nucli d'urani en dos punts, A i B, situats a 0,008~nm i 0,008~pm, respectivament, del centre d'aquest nucli.
\end{apartat}

\begin{apartat}{1}
Quanta energia cinètica ha de tenir, com a mínim, la partícula $\alpha$ quan passa pel punt A per a arribar fins al punt B? (Ignoreu la influència que els electrons pròxims puguin tenir.)
\end{apartat}

\dades{%
  $k = \dfrac{1}{4\pi\varepsilon_0} = 8,99\times 10^{9}\ \mathrm{N\,m^2\,C^{-2}}$.\\
  Càrrega elemental $= 1,60\times 10^{-19}\ \mathrm{C}$.\\
  Nombre atòmic de l'urani $= 92$.}
